% ============================================================
% sec/04stage3.tex - Etapa 3 (07/11 - semana 16)
% ============================================================
\section{Etapa 3: agente, validaci\'on, \'etica y art\'iculo final}
\label{sec:stage3}

\notebox[Objetivo de la etapa]{Convertir el modelo final de la Etapa 2 en un
componente que una persona pueda usar a trav\'es de un agente, validar ese
agente con un protocolo escrito antes de ejecutar, declarar la \'etica de su
despliegue y cerrar el art\'iculo. El modelo no se reentrena ni se recalibra.
Peso: 25\,\%.}

\subsection{Resumen de ramas}

\begin{table}[H]
  \centering
  \caption{Ramas de la Etapa 3 y su resultado.}
  \label{tab:stage3summary}
  \footnotesize
  \begin{tabularx}{\textwidth}{@{}P{5.0cm} P{1.6cm} P{2.6cm} Y@{}}
    \toprule
    \textbf{Rama} & \textbf{Fechas} & \textbf{Responsables} & \textbf{Resultado}\\
    \midrule
    \code{07MultiAgentIntegration} & 07/11 & Adriel & Herramientas del modelo, cuatro agentes, registro, servidor MCP, orquestador y subagentes de Claude Code (PR \#13)\\
    {\scriptsize\code{08AgentValidationAndSimulator}} & 08/11 & Adriel; validaci\'on humana de los tres & Validaci\'on N1 y N2, an\'alisis de fragilidad y simulador gr\'afico (PR \#14)\\
    {\scriptsize\code{09EthicsPaperAndDocumentation}} & 08/11 & Adriel; revisi\'on de Daniel y Sebasti\'an & \'Etica del despliegue, art\'iculo final, documentaci\'on y uso de IA\\
    \code{10DefenseAndShowcase} & Semana 16 & Equipo & Pendiente: defensa, video y entrega\\
    \bottomrule
  \end{tabularx}
\end{table}

\subsection{Lo que recibi\'o de la Etapa 2}

\begin{table}[H]
  \centering
  \caption{Piezas que la Etapa 2 dej\'o listas para el agente.}
  \label{tab:entregaetapa3}
  \footnotesize
  \begin{tabularx}{\textwidth}{@{}P{4.2cm} Y@{}}
    \toprule
    \textbf{Pieza} & \textbf{Detalle}\\
    \midrule
    Modelo final & B4 entrenado en el escenario A, sin calibrador. En Kaggle, \code{modelo\_\allowbreak{}final/\allowbreak{}b4\_\allowbreak{}victoria.joblib} con su ficha (origen, SHA-256, versi\'on de scikit-learn y m\'etricas)\\
    Interfaz & \code{ptcg\_\allowbreak{}wp.predictor.PredictorVictoria}: recibe la observaci\'on del jugador que act\'ua y devuelve la probabilidad de que gane. En preparaci\'on y en estados ya resueltos devuelve \code{None}\\
    Garant\'ia de consistencia & \code{tests/\allowbreak{}test\_\allowbreak{}predictor.py} comprueba que la observaci\'on produce las mismas caracter\'isticas, en el mismo orden, que la tabla de entrenamiento\\
    Explicaciones & SHAP mostr\'o que el modelo decide por premios, descartes y mazo; las explicaciones del agente se dan en esos t\'erminos\\
    \bottomrule
  \end{tabularx}
\end{table}

Tambi\'en recibi\'o tres condiciones de uso: desconfiar de las probabilidades
extremas al inicio de la partida, comprobar que el agente no sobrevalore
dejar al rival sin mazo y no confiar en la calibraci\'on con datos de otro
periodo sin volver a medirla. Las tres quedaron como advertencias del agente
o como l\'imites declarados.

\subsection{C\'omo se construy\'o, parte por parte}

\begin{figure}[H]
  \centering
  \begin{tikzpicture}[
    font=\footnotesize,
    caja/.style={draw=tecsteel, fill=boxbg, rounded corners=2pt,
                 text width=3.0cm, align=center, minimum height=1.0cm},
    rama/.style={font=\scriptsize\itshape, text=labelgray},
    flecha/.style={-{Stealth[length=2mm]}, thick, tecsteel}]
    \node[caja] (m) {Modelo final\\(Etapa 2)};
    \node[caja, right=0.6cm of m] (h) {Herramientas\\y verificador};
    \node[caja, right=0.6cm of h] (a) {Agentes, MCP\\y orquestador};
    \node[caja, below=1.0cm of a] (v) {Validaci\'on\\N1 y N2};
    \node[caja, left=0.6cm of v] (s) {Simulador y\\validaci\'on humana};
    \node[caja, left=0.6cm of s] (e) {\'Etica, art\'iculo\\y documentaci\'on};
    \draw[flecha] (m) -- (h);
    \draw[flecha] (h) -- (a);
    \draw[flecha] (a) -- (v);
    \draw[flecha] (v) -- (s);
    \draw[flecha] (s) -- (e);
    \node[rama, above=0.05cm of h] {rama 07};
    \node[rama, below=0.05cm of s] {rama 08};
    \node[rama, below=0.05cm of e] {rama 09};
  \end{tikzpicture}
  \caption{Orden de construcci\'on de la Etapa 3. Cada pieza se prob\'o antes
  de empezar la siguiente.}
  \label{fig:flujo3}
\end{figure}

% ============================================================
\subsection{Rama 07: integraci\'on del agente}
\label{subsec:rama07}

\jentry{2026-11-07}{Herramientas, agentes y servidor MCP}%
       {Adriel}{07MultiAgentIntegration}

Se escribi\'o primero la arquitectura (\code{07-11-26\_\allowbreak{}ArquitecturaDelAgente.md})
y despu\'es el c\'odigo en \code{src/\allowbreak{}ptcg\_\allowbreak{}agente/}, separado del
c\'odigo de entrenamiento. Hay seis herramientas del modelo (estimar,
explicar, evaluar una partida, comparar estados, revisar el dominio y
describir el modelo), cuatro agentes y un registro JSONL de cada
invocaci\'on. El agente se usa desde Claude Code o Claude Desktop, con el
servidor MCP \code{ptcg-victoria}, o desde un orquestador en Python.

\decision{Usar la membres\'ia de Claude del equipo, sin clave de API: Claude
Code y Claude Desktop con un servidor MCP, y Claude Code sin interfaz
(\code{claude -p}) para el planificador y el redactor del
orquestador.}{Una clave de API ser\'ia otra credencial que proteger y otro
costo. Con la membres\'ia, cada integrante usa su propia cuenta.}{Una clave
de API de Anthropic o un modelo local m\'as d\'ebil.}

\decision{El planificador y el redactor son dos llamadas separadas a Claude,
sin herramientas; el ejecutor y el verificador son c\'odigo.}{As\'i el LLM
nunca calcula una probabilidad: solo decide qu\'e herramientas llamar y
redacta con resultados ya aprobados. El verificador revisa que cada cifra de
la respuesta salga de una herramienta, que aparezcan las advertencias
obligatorias y que no se prometa certeza.}{Un solo agente con todas las
herramientas, que mezcla decidir, calcular y redactar en la misma llamada.}

\milestone{El agente da la misma probabilidad que el modelo final en 642
estados de 4 partidas reales (diferencia m\'axima 0.0).}{Parte 3 del cuaderno
\code{07\_\allowbreak{}agente}; sesi\'on de Claude Code en
\code{08-11-26\_\allowbreak{}RegistroSesionClaudeCode.jsonl} y \code{.png}, donde el
verificador rechaz\'o un primer borrador y aprob\'o el corregido.}

\aiuse{Claude}{proponer la arquitectura e implementar las herramientas, los
agentes, el servidor MCP y sus pruebas}{Adriel ejecut\'o el cuaderno 07,
instal\'o el agente en su PC y revis\'o cada respuesta guardada como evidencia}

% ============================================================
\subsection{Rama 08: validaci\'on y simulador}
\label{subsec:rama08}

\jentry{2026-11-08}{Protocolo y validaci\'on en dos niveles}%
       {Adriel}{08AgentValidationAndSimulator}

El protocolo (\code{08-11-26\_\allowbreak{}ProtocoloValidacionAgente.md}) fij\'o,
antes de ejecutar, los escenarios, el comportamiento esperado y el criterio
de aprobaci\'on de cada caso. La validaci\'on tiene dos niveles: N1 prueba
las capas deterministas sin LLM (cuaderno 09, en Colab) y N2 prueba el
agente completo con Claude (cuaderno 10, en la PC).

\blocker{Kaggle respond\'ia 429 (demasiadas solicitudes) al bajar en paralelo
las 166 repeticiones de los casos.}{La parte 1 del cuaderno 09 fallaba a
mitad de la descarga.}{Descargas una por una, con espera y reintento, en una
carpeta que conserva lo ya bajado.}

\blocker{Kaggle respond\'ia 404 al bajar \code{partidas.json.gz}.}{El cuaderno
10 no encontraba las partidas de validaci\'on.}{Guardar \code{partidas.json}
sin comprimir y borrar el archivo anterior del dataset.}

\blocker{La herramienta \code{kaggle} 2.0.2 de Colab respond\'ia 400 al crear
una versi\'on, sin decir el motivo.}{No se pod\'ia publicar el cat\'alogo de
cartas del simulador.}{Fijar \code{kaggle>=2.2.5} en \code{requirements.txt}
y, si vuelve a ocurrir, repetir con la API de Python para mostrar el motivo
sin imprimir el token.}

La primera corrida de N1 dio 32 de 35 casos aprobados, 2 fallidos y 1 de
registro. Los dos fallos eran de presentaci\'on, no del modelo:
\begin{itemize}
  \item \textbf{L-09.} La partida m\'as larga (1398 pasos) tardaba 56.6\,s
        porque se llamaba al modelo una vez por paso. Ahora se predice toda
        la partida en una llamada: 0.9\,s, con las mismas probabilidades.
  \item \textbf{L-10.} Sin Pok\'emon activo, la explicaci\'on mostraba el
        valor interno $-1$ como si fueran puntos de vida. Ahora dice ``no hay
        Pok\'emon activo''.
\end{itemize}

\decision{Reportar la corrida original junto a la corregida, sin cambiar los
criterios del protocolo.}{Mostrar solo la corrida corregida ocultar\'ia que
la validaci\'on s\'i encontr\'o defectos reales.}{Ajustar los umbrales de
L-09 o borrar la primera corrida.}

\milestone{N1 corregido: 34 de 35 casos aprobados y 1 de registro. N2: 31 de
31 consultas aprobadas, herramienta correcta en 18 de 18, cifras trazables en
31 de 31, 28 de 31 aprobadas en la primera ronda del redactor, 0 respuestas en
la plantilla, latencia p95 de 31\,s y costo equivalente de 1.72\,USD.}{%
\code{08-11-26\_\allowbreak{}ResultadosValidacionAgente.md},
\code{08-11-26\_\allowbreak{}ValidacionN1Corregida.csv} y
\code{08-11-26\_\allowbreak{}ValidacionN2.csv}.}

El hallazgo m\'as importante fue de fragilidad: 50 de 60 errores muy seguros
de medio y final de la partida llegan sin advertencia, y el agente los
presenta como fiables. El verificador detecta cifras inventadas y
advertencias omitidas, pero no puede saber cu\'ando el modelo se equivoca con
seguridad.

\jentry{2026-11-08}{Simulador gr\'afico y validaci\'on humana}%
       {Adriel; los tres en la validaci\'on humana}%
       {08AgentValidationAndSimulator}

El simulador (\code{python -m ptcg\_\allowbreak{}agente.simulador}) es una p\'agina web
local que muestra el tablero de los dos jugadores, deja editar el estado y
recalcula la probabilidad, las contribuciones SHAP y las advertencias en cada
cambio. Usa solo la biblioteca est\'andar de Python y JavaScript sin
librer\'ias. El agente lo abre con la herramienta \code{abrir\_\allowbreak{}simulador}.

\decision{Jugadores fijos en pantalla (el 1 arriba, el 0 abajo) y la cifra
grande siempre del jugador~0.}{En la primera versi\'on la cifra era del
jugador que actuaba, as\'i que cambiaba de due\~no al avanzar la partida y
confund\'ia al usuario.}{Mostrar la probabilidad del jugador que act\'ua, que es
la salida directa del modelo.}

\decision{Agregar un ``estado nuevo'' vac\'io (turno 0, sin cartas ni
Pok\'emon) y un cat\'alogo con los PV m\'aximos de las 362 cartas vistas como
activo.}{Permite armar un estado desde cero, sin depender de una partida, y
elegir cualquier carta con sus PV reales.}{Solo editar estados de partidas
reales.}

\decision{En la validaci\'on humana, 5 pruebas distintas por integrante (15 en
total) con la direcci\'on esperada fijada antes.}{Una primera versi\'on
repet\'ia las mismas pruebas para los tres y no a\~nad\'ia
informaci\'on.}{Tres sesiones id\'enticas.}

\milestone{El simulador funcion\'o en las tres computadoras y 10 de 15 pruebas
tuvieron la direcci\'on esperada: todas las de premios, mazo, PV y banca. Las
condiciones especiales casi no movieron la probabilidad, y m\'as energ\'ias en
el activo la bajaron (asociaci\'on, no causa).}{%
\code{08-11-26\_\allowbreak{}ValidacionHumanaSimulador.md} y las capturas
\code{08-11-26\_\allowbreak{}Simulador*.png}.}

\aiuse{Claude}{implementar la bater\'ia, los cuadernos 09 y 10 y el simulador,
y diagnosticar los errores de Kaggle}{Adriel ejecut\'o N1 y N2 y revis\'o a mano
Q-01, Q-02 y Q-07; cada integrante hizo sus propias pruebas con el simulador}

% ============================================================
\subsection{Rama 09: \'etica, art\'iculo y documentaci\'on}
\label{subsec:rama09}

\jentry{2026-11-08}{\'Etica del despliegue y art\'iculo final}%
       {Adriel; revisi\'on de Daniel y Sebasti\'an}{09EthicsPaperAndDocumentation}

La \'etica (\code{08-11-26\_\allowbreak{}EticaDelDespliegue.md}) declara 8 riesgos
con su gravedad y probabilidad, los datos sensibles, una mitigaci\'on con su
archivo y su prueba para cada riesgo, 9 l\'imites y el uso aceptable. El
art\'iculo sum\'o la secci\'on del agente (arquitectura, protocolo, resultados
y fragilidad), con la consulta real en Claude Code y el simulador como
figuras; la \'etica se reescribi\'o con foco en el despliegue; y la discusi\'on
y las conclusiones incluyen el agente y el objetivo~7.

\decision{No ofrecer el sistema como servicio p\'ublico: todo corre en la
computadora del usuario.}{El sistema no puede impedir por s\'i mismo el uso en
partidas en curso o en apuestas; sin un servicio p\'ublico, ese uso es menos
f\'acil y el uso aceptable queda declarado.}{Publicar el simulador en la
web.}

\decision{Declarar como l\'imite, y no ocultar, que la mayor\'ia de los errores
muy seguros no traen advertencia.}{Es el riesgo m\'as grave (confianza
excesiva) y el lector debe conocerlo. La mejora propuesta, mostrar con qu\'e
frecuencia fallan las predicciones muy seguras, queda como trabajo
futuro.}{Agregar una advertencia nueva sin validarla con un protocolo.}

\milestone{Art\'iculo final de 29 p\'aginas que compila sin errores, con los
siete objetivos cumplidos; documentaci\'on del repositorio actualizada para
ejecutar el modelo, el agente y el simulador.}{Commits de la rama
\code{09EthicsPaperAndDocumentation}.}

Daniel y Sebasti\'an revisaron, con apoyo de Claude, que el desarrollo de la
etapa siguiera el plan y las reglas del equipo y que todo funcionara bien
seg\'un su propia evaluaci\'on humana, adem\'as de hacer sus pruebas de la
validaci\'on humana (\code{01UsoDeIA/\allowbreak{}Etapa3UsoDeIA.md}).

% ============================================================
\subsection{C\'omo aseguramos que el agente fuera confiable}
\label{subsec:verificacion3}

\begin{itemize}
  \item \textbf{Mismo modelo, comprobado.} El agente solo arranca si el modelo
        y sus archivos coinciden por SHA-256 con los publicados, y con la
        misma versi\'on de scikit-learn.
  \item \textbf{Ninguna cifra inventada.} Toda probabilidad sale de una
        herramienta y queda en el registro; el verificador rechaza las
        respuestas con cifras que no salieron de una herramienta.
  \item \textbf{Sin informaci\'on oculta.} Ni el modelo ni el agente usan la
        mano del rival ni el contenido de los premios; una prueba lo comprueba.
  \item \textbf{Pruebas autom\'aticas.} Pasaron de 109 al cerrar la Etapa 2 a
        145, con las del agente, el orquestador (con un LLM de prueba), el
        servidor MCP, la validaci\'on y el simulador. Corren como segunda
        celda de los cuadernos 07 a 10.
  \item \textbf{Protocolo antes de ejecutar.} Los criterios de aprobaci\'on se
        fijaron por escrito y no cambiaron al ver los resultados.
  \item \textbf{Evaluaci\'on humana.} Cada integrante prob\'o el simulador en su
        computadora y revis\'o el funcionamiento por su cuenta.
\end{itemize}

% ============================================================
\subsection{Retrospectiva}

\begin{table}[H]
  \centering
  \caption{Retrospectiva de la Etapa 3.}
  \label{tab:stage3retro}
  \footnotesize
  \begin{tabularx}{\textwidth}{@{}P{4.0cm} Y@{}}
    \toprule
    \textbf{Pregunta} & \textbf{Respuesta}\\
    \midrule
    ¿Qu\'e funcion\'o? & Separar decidir y redactar (LLM) de calcular y verificar (c\'odigo). Partir de un predictor ya probado. Escribir el protocolo antes de ejecutar. Mantener Kaggle como almac\'en \'unico tambi\'en para los archivos del agente.\\
    ¿Qu\'e nos bloque\'o? & Los l\'imites y cambios de Kaggle (429, archivos \code{.gz} y el error 400 de la herramienta vieja) y la espera mientras Kaggle procesa cada versi\'on.\\
    ¿Qu\'e estimamos mal? & El costo de evaluar una partida paso a paso (L-09) y la claridad del simulador: la primera versi\'on, con la cifra del jugador que act\'ua, confund\'ia.\\
    ¿Qu\'e queda abierto? & Advertir cu\'ando una predicci\'on muy segura suele fallar, que el modelo use mejor las condiciones especiales y la defensa con el video demostrativo (rama 10).\\
    \bottomrule
  \end{tabularx}
\end{table}

\subsection{Trabajo pendiente heredado}

\begin{itemize}
  \item Repetir los experimentos principales con varias semillas.
  \item Agregar una caracter\'istica con los premios que entrega cada
        Pok\'emon activo, que hoy el modelo aproxima con sus puntos de vida
        m\'aximos.
  \item Entrenar la variante de 64 caracter\'isticas (estado completo m\'as
        identificadores de carta).
\end{itemize}